% !TEX root = 00_main.tex
% =================================================================
%  マクロ定義
% =================================================================

% 図表・章・式の参照
%   ラベルは \label{fig:xxx} \label{tab:xxx} \label{chap:xxx} \label{eq:xxx} の形で付ける．
%   本文中では \figref{xxx} → 「図 1.1」のように出力される．
\newcommand\figref[1]{図~\ref{fig:#1}}
\newcommand\tabref[1]{表~\ref{tab:#1}}
\newcommand\chapref[1]{第~\ref{chap:#1}~章}
\newcommand\secref[1]{\ref{sec:#1}~節}
\newcommand\equref[1]{式~(\ref{eq:#1})}

\renewcommand{\bibname}{参考文献}

% 行間
\def\baselinestretch{1.15} \small\normalsize

% 執筆中のメモ（赤字で表示）
%   \todo{あとで図を追加} のように使う．
%   提出前に下の \showtodotrue を \showtodofalse にすると全て非表示になる．
\newif\ifshowtodo
\showtodotrue
\newcommand{\todo}[1]{\ifshowtodo\textcolor{red}{\textbf{[TODO: #1]}}\fi}
